% GENERATED FILE. Edit canonical lessons, then run npm run generate:books.
% canonical-source-hash: fnv1a64:69e99605c2dac4d4
% canonical-lessons: KA-C41-this, KA-C41-that, KA-C41-here, KA-C41-there, KA-C41-who, KA-C41-where, KA-C41-deixis-system

\chapter{Pointing, and Asking}
\label{ch:kn-pointing-and-asking}
\hlchaptermodality{41}

\begin{chapteropening}
\textbf{By the end of this chapter:} \emph{I can point at something in Kannada and ask who or where it is — and I can say what the six words have in common.}

Say the near word, the far word and the question word in a row, and name the part that changed.
\end{chapteropening}

\section[idu]{\kn{ಇದು} (idu) — this one — the thing near me}

\label{lesson:KA-C41-this}



\begin{quote}
You can already name things. This is how you POINT at them.
\end{quote}

\subsection*{You'll want to know: \kn{ಇದು}}
\textbf{\kn{ಇದು}} — \emph{idu} — this one — the thing near me.

Say it, and point while you say it. That is the whole word: it does not mean anything on its own, it means whatever your finger is on.

Its partner is \textbf{\kn{ಅದು}} \emph{adu}, which you will meet in a moment. The two of them differ by one sound at the front — \textbf{i-} for near — and that is not a coincidence. It is the whole system, and the last lesson of this chapter says so.

\subsection*{Your turn}
Say these aloud:

\begin{itemize}
\raggedright
  \item \textquotedblleft{}\kn{ಇದು}\textquotedblright{} three times, pointing at something different each time
  \item it once with a word you already know after it
\end{itemize}

\begin{tcolorbox}[breakable,colback=teal!4,colframe=teal!35!black,title={Before you move on}]
What is \emph{idu}? (this one — the thing near me.) What sound does it start with, and does that sound mean near, far, or a question? (\textbf{i-} — near.)
\end{tcolorbox}
\section[adu]{\kn{ಅದು} (adu) — that one — the thing over there}

\label{lesson:KA-C41-that}



\begin{quote}
You can already name things. This is how you POINT at them.
\end{quote}

\subsection*{You'll want to know: \kn{ಅದು}}
\textbf{\kn{ಅದು}} — \emph{adu} — that one — the thing over there.

Say it, and point while you say it. That is the whole word: it does not mean anything on its own, it means whatever your finger is on.

Its partner is \textbf{\kn{ಇದು}} \emph{idu}, which you will meet just now. The two of them differ by one sound at the front — \textbf{a-} for far — and that is not a coincidence. It is the whole system, and the last lesson of this chapter says so.

\subsection*{Your turn}
Say these aloud:

\begin{itemize}
\raggedright
  \item \textquotedblleft{}\kn{ಅದು}\textquotedblright{} three times, pointing at something different each time
  \item it once with a word you already know after it
\end{itemize}

\begin{tcolorbox}[breakable,colback=teal!4,colframe=teal!35!black,title={Before you move on}]
What is \emph{adu}? (that one — the thing over there.) What sound does it start with, and does that sound mean near, far, or a question? (\textbf{a-} — far.)
\end{tcolorbox}
\section[illi]{\kn{ಇಲ್ಲಿ} (illi) — here — where I am}

\label{lesson:KA-C41-here}



\begin{quote}
You can already name things. This is how you POINT at them.
\end{quote}

\subsection*{You'll want to know: \kn{ಇಲ್ಲಿ}}
\textbf{\kn{ಇಲ್ಲಿ}} — \emph{illi} — here — where I am.

Say it, and point while you say it. That is the whole word: it does not mean anything on its own, it means whatever your finger is on.

Its partner is \textbf{\kn{ಅಲ್ಲಿ}} \emph{alli}, which you will meet in a moment. The two of them differ by one sound at the front — \textbf{i-} for near — and that is not a coincidence. It is the whole system, and the last lesson of this chapter says so.

\subsection*{Your turn}
Say these aloud:

\begin{itemize}
\raggedright
  \item \textquotedblleft{}\kn{ಇಲ್ಲಿ}\textquotedblright{} three times, pointing at something different each time
  \item it once with a word you already know after it
\end{itemize}

\begin{tcolorbox}[breakable,colback=teal!4,colframe=teal!35!black,title={Before you move on}]
What is \emph{illi}? (here — where I am.) What sound does it start with, and does that sound mean near, far, or a question? (\textbf{i-} — near.)
\end{tcolorbox}
\section[alli]{\kn{ಅಲ್ಲಿ} (alli) — there — where I am not}

\label{lesson:KA-C41-there}



\begin{quote}
You can already name things. This is how you POINT at them.
\end{quote}

\subsection*{You'll want to know: \kn{ಅಲ್ಲಿ}}
\textbf{\kn{ಅಲ್ಲಿ}} — \emph{alli} — there — where I am not.

Say it, and point while you say it. That is the whole word: it does not mean anything on its own, it means whatever your finger is on.

Its partner is \textbf{\kn{ಇಲ್ಲಿ}} \emph{illi}, which you will meet just now. The two of them differ by one sound at the front — \textbf{a-} for far — and that is not a coincidence. It is the whole system, and the last lesson of this chapter says so.

\subsection*{Your turn}
Say these aloud:

\begin{itemize}
\raggedright
  \item \textquotedblleft{}\kn{ಅಲ್ಲಿ}\textquotedblright{} three times, pointing at something different each time
  \item it once with a word you already know after it
\end{itemize}

\begin{tcolorbox}[breakable,colback=teal!4,colframe=teal!35!black,title={Before you move on}]
What is \emph{alli}? (there — where I am not.) What sound does it start with, and does that sound mean near, far, or a question? (\textbf{a-} — far.)
\end{tcolorbox}
\section[yāru]{\kn{ಯಾರು} (yāru) — who? — asking about a person}

\label{lesson:KA-C41-who}



\begin{quote}
You can already name things. This is how you POINT at them.
\end{quote}

\subsection*{You'll want to know: \kn{ಯಾರು}}
\textbf{\kn{ಯಾರು}} — \emph{yāru} — who? — asking about a person.

Say it, and point while you say it. That is the whole word: it does not mean anything on its own, it means whatever your finger is on.

Its partner is \textbf{\kn{ಎಲ್ಲಿ}} \emph{elli}, which you will meet in a moment. The two of them differ by one sound at the front — \textbf{e-} for ask — and that is not a coincidence. It is the whole system, and the last lesson of this chapter says so.

\subsection*{Your turn}
Say these aloud:

\begin{itemize}
\raggedright
  \item \textquotedblleft{}\kn{ಯಾರು}\textquotedblright{} three times, pointing at something different each time
  \item it once with a word you already know after it
\end{itemize}

\begin{tcolorbox}[breakable,colback=teal!4,colframe=teal!35!black,title={Before you move on}]
What is \emph{yāru}? (who? — asking about a person.) What sound does it start with, and does that sound mean near, far, or a question? (\textbf{e-} — ask.)
\end{tcolorbox}
\section[elli]{\kn{ಎಲ್ಲಿ} (elli) — where? — asking about a place}

\label{lesson:KA-C41-where}



\begin{quote}
You can already name things. This is how you POINT at them.
\end{quote}

\subsection*{You'll want to know: \kn{ಎಲ್ಲಿ}}
\textbf{\kn{ಎಲ್ಲಿ}} — \emph{elli} — where? — asking about a place.

Say it, and point while you say it. That is the whole word: it does not mean anything on its own, it means whatever your finger is on.

Its partner is \textbf{\kn{ಯಾರು}} \emph{yāru}, which you will meet just now. The two of them differ by one sound at the front — \textbf{e-} for ask — and that is not a coincidence. It is the whole system, and the last lesson of this chapter says so.

\subsection*{Your turn}
Say these aloud:

\begin{itemize}
\raggedright
  \item \textquotedblleft{}\kn{ಎಲ್ಲಿ}\textquotedblright{} three times, pointing at something different each time
  \item it once with a word you already know after it
\end{itemize}

\begin{tcolorbox}[breakable,colback=teal!4,colframe=teal!35!black,title={Before you move on}]
What is \emph{elli}? (where? — asking about a place.) What sound does it start with, and does that sound mean near, far, or a question? (\textbf{e-} — ask.)
\end{tcolorbox}
\section[i- / a- / e-]{i- / a- / e- — six words, one machine}

\label{lesson:KA-C41-deixis-system}



\begin{quote}
Say the six words from this chapter out loud, in order. Listen to how they begin.
\end{quote}

\begin{grammarlens}[title={Grammar lens: one slot changes}]
You did not learn six words. You learned \textbf{three beginnings} and reused them.

\noindent
\begin{tabularx}{\linewidth}{@{}>{\raggedright\arraybackslash}X>{\raggedright\arraybackslash}X@{}}
\toprule
\textbf{meaning} & \textbf{word} \\
\midrule
near & \textbf{\kn{ಇಲ್ಲಿ}} \emph{illi} \\
far & \textbf{\kn{ಅಲ್ಲಿ}} \emph{alli} \\
question & \textbf{\kn{ಎಲ್ಲಿ}} \emph{elli} \\
\bottomrule
\end{tabularx}

Kannada's place words are the shortest of the four: \kn{ಇಲ್ಲಿ}, \kn{ಅಲ್ಲಿ}, \kn{ಎಲ್ಲಿ}.

Change the front, and the meaning walks from \textbf{near} to \textbf{far} to \textbf{a question}. Nothing else in the word moves.

The same three beginnings run the thing-words too: \textbf{\kn{ಇದು}} \emph{idu} and \textbf{\kn{ಅದು}} \emph{adu} start exactly like the first two rows. This chapter did not teach you the matching question word for things — that one arrives later, and when it does you will already know what it has to start with.

This is worth more than the six words are. When you meet a new word in this family later, you will not have to be taught all three of it — you will be taught one and work out the others.
\end{grammarlens}

\subsection*{Your turn}
\begin{itemize}
\raggedright
  \item \emph{Say it:} the near word, then the far word, then the question word
  \item \emph{Say it:} them again for place, not for things
  \item \emph{Notice:} the part that did NOT change
\end{itemize}

\begin{tcolorbox}[breakable,colback=teal!4,colframe=teal!35!black,title={Before you move on}]
Which beginning means near? (\textbf{i-}) Far? (\textbf{a-}) And which one turns the word into a question? (\textbf{e-})
\end{tcolorbox}
